% Generated draft: requires mathematical and source review.
\documentclass[11pt]{article}
\usepackage[margin=25mm]{geometry}
\usepackage{fontspec}
\setmainfont{Latin Modern Roman}
\usepackage{amsmath,amssymb,mathtools,unicode-math}
\setmathfont{Latin Modern Math}
\usepackage{enumitem,needspace,xurl,hyperref}
\setlength{\parindent}{0pt}
\setlength{\parskip}{5pt}
\setlength{\emergencystretch}{3em}
\newcommand{\N}{\mathbb{N}}
\newcommand{\Z}{\mathbb{Z}}
\newcommand{\Q}{\mathbb{Q}}
\newcommand{\R}{\mathbb{R}}
\newcommand{\C}{\mathbb{C}}
\newcommand{\F}{\mathbb{F}}
\newcommand{\A}{\mathbb{A}}
\newcommand{\PP}{\mathbb{P}}
\newcommand{\E}{\mathbb{E}}
\newcommand{\eps}{\varepsilon}
\newcommand{\dd}{\,\mathrm{d}}
\begin{document}
\section*{UnsolvedMath ID: 20002349}
\textbf{Source problem code:} catalogue entry
\subsection{Definitions and mathematical statement}
A subset of \(\Q^k\) is Diophantine over \(\Q\) if it is defined by an existential system of polynomial equations with rational coefficients. A Diophantine interpretation of \(\Z\) in \(\Q\), allowing a quotient, consists of a Diophantine domain \(S\subseteq\Q^k\), a Diophantine equivalence relation \(E\) on \(S\), and Diophantine relations inducing addition, multiplication, and the distinguished constants on \(S/E\), together with an isomorphism of the resulting quotient ring with \(\Z\). Write \(\pi:S\to\Z\) for the corresponding surjection, so \(E(x,y)\) means \(\pi(x)=\pi(y)\). The map \(\pi\) need not be a definable choice of representatives.

Mazur's finite-components conjecture, in the form relevant here, says that for every algebraic variety \(V\) over \(\Q\), the closure of \(V(\Q)\) in \(V(\R)\), with its usual real topology, has finitely many connected components. Cornelissen's Question 12 asks whether
\[
 \Z\text{ has a Diophantine quotient interpretation in }\Q
 \quad\Longrightarrow\quad
 \text{Mazur's finite-components conjecture fails}.
\]
No uniform finite bound or finiteness of the \(E\)-classes is assumed in this original implication.

The supplied partial result adds a uniform bound: if there is one integer \(M\ge1\) such that
\[
 \forall n\in\Z\qquad 1\le|\pi^{-1}(n)|\le M,
\]
then Mazur's conjecture is false. Equivalently, under Mazur's conjecture no such interpretation can have all class sizes bounded by one finite constant. This does not exclude interpretations with unbounded finite classes or with infinite classes.

The aid describes a selector-free construction for this bounded-class assertion. For a finite class \(C\subseteq\Q^k\) of size \(r\), associate the polynomial
\[
 P_C(T_0,\ldots,T_k)
 =\prod_{x=(x_1,\ldots,x_k)\in C}
 (1+T_0+x_1T_1+\cdots+x_kT_k).
\]
Use its coefficient vector, in a fixed monomial order through degree \(r\), as a whole-class invariant. The normalization fixes the scalar ambiguity in its linear factors, so distinct classes have distinct coefficient vectors; no ordering or distinguished member of a class is required. The aid claims that one can choose an infinite recursively enumerable set of integer labels whose classes all have one fixed size \(r\), obtain a Diophantine set of these coefficient vectors, and make every recursively enumerable subfamily Diophantine by interpreting integer Diophantine conditions. It then applies Cornelissen--Zahidi topological thinning to obtain a contradiction to the finite-components conjecture.

The exact-size restriction is a material part of that argument: merely saying that at least \(r\) distinct equivalent points exist does not say that they exhaust a larger class, and the absence of a further point is not automatically an existential polynomial condition. The supplied summary asserts a valid stratum extraction but gives no detailed proof. Thus the bounded-class assertion and its Diophantine coding steps are retained as attributed partial-result claims requiring mathematical verification.

The aid also reports a topological example showing that discreteness of an abstract quotient alone does not force disconnected real closure of its representatives. In particular, no topology or discreteness conclusion about \(S\subseteq\Q^k\) follows merely by identifying \(S/E\) abstractly with \(\Z\). This is why the general question is not replaced by the bounded-class result. The trailing page number and workshop editor header are extraction artifacts, not mathematical hypotheses.
\subsection{Short English statement}
Does a Diophantine quotient interpretation of \(\Z\) in \(\Q\) contradict Mazur's finite-components conjecture? The supplied aid claims this under a uniform finite bound on equivalence-class sizes, using whole-class coefficient invariants, while leaving the unrestricted quotient question open.
\subsection{Sources}
\begin{itemize}
\item Local registry, permanent UnsolvedMath ID \(20002349\), exact public code \texttt{catalogue entry}; canonical location \texttt{aim-logic-notes.json}, \texttt{notes[124]}.
\item AIM, \emph{Problems related to ``Extensions of Hilbert's tenth problem''}, Question 12 attributed to Cornelissen, official page 3, PDF page index 2 according to the supplied verification field. Workshop moderated by B. Poonen and T. Scanlon, notes by J. Demeyer; supplied reference \url{https://aimath.org/WWN/hilberts10th/hilberts10th.pdf}. No source inspection was performed here.
\item Cornelissen--Zahidi, \emph{Contemporary Mathematics} \textbf{270}, 253--260, as cited in the source. No title, year, or theorem number is supplied. The aid invokes their topological thinning argument.
\item Supplied research attempt 1: \texttt{partial\_result}, run status \texttt{open}, rigor label \texttt{proved}, low novelty confidence; registry status \texttt{partially\_solved}, classification \texttt{PARTIAL-PROGRESS}, created and updated 14 August 2026. The bounded-class extension and topological example are machine-generated aid claims, not peer reviewed; the aid requires independent verification before citation.
\end{itemize}
\end{document}
